\chapter{Error Handling: Dealing with Mistakes Gracefully}

\section{Introduction: Real-Life Hook}

Imagine a student writing an exam.

If the student makes a small mistake:
\begin{itemize}
    \item The teacher does not stop the exam
    \item The student continues answering other questions
\end{itemize}

But if a program makes a mistake:
\begin{itemize}
    \item It suddenly stops
    \item It shows a scary error message
\end{itemize}

Can we teach programs to handle mistakes calmly?

Yes — using \textbf{Error Handling}.

\section{What Is an Error?}

\begin{definition}{Error}
An error is a problem in a program that stops normal execution.
\end{definition}

Examples:
\begin{itemize}
    \item Dividing by zero
    \item Opening a missing file
    \item Writing wrong data type
\end{itemize}

\section{Why Error Handling Is Important}

Without error handling:
\begin{itemize}
    \item Program crashes
    \item User feels confused
\end{itemize}

With error handling:
\begin{itemize}
    \item Program continues safely
    \item User gets friendly messages
\end{itemize}

\section{Types of Errors (Brief)}

\begin{itemize}
    \item Syntax errors
    \item Runtime errors
    \item Logical errors
\end{itemize}

This chapter focuses on \textbf{runtime errors}.

\section{Try and Except Blocks}

Python uses \texttt{try-except} to handle errors.

\begin{verbatim}
try:
    number = int(input("Enter a number: "))
    print(10 / number)
except:
    print("Something went wrong")
\end{verbatim}

If error occurs, program jumps to \texttt{except}.

\section{Handling Specific Errors}

\begin{verbatim}
try:
    x = int(input("Enter number: "))
    y = 10 / x
except ZeroDivisionError:
    print("Cannot divide by zero")
except ValueError:
    print("Please enter a valid number")
\end{verbatim}

\section{Using Else Block}

\begin{verbatim}
try:
    num = int(input("Enter number: "))
    print(100 / num)
except:
    print("Error occurred")
else:
    print("Calculation successful")
\end{verbatim}

Else runs when no error occurs.

\section{Finally Block}

\begin{verbatim}
try:
    file = open("data.txt", "r")
    print(file.read())
except:
    print("File not found")
finally:
    print("Program finished")
\end{verbatim}

Finally always runs.

\section{Real-Life Example: ATM Machine}

ATM behavior:
\begin{itemize}
    \item Wrong PIN $\rightarrow$ shows message
    \item Correct PIN $\rightarrow$ allows transaction
    \item Machine never crashes
\end{itemize}

Programs should behave like ATMs.

\section{Mini-Project: Safe Calculator}

\begin{verbatim}
try:
    a = int(input("Enter first number: "))
    b = int(input("Enter second number: "))
    print(a / b)
except ZeroDivisionError:
    print("Division by zero not allowed")
except ValueError:
    print("Invalid input")
else:
    print("Calculation completed")
\end{verbatim}

Students learn:
\begin{itemize}
    \item Safe input handling
    \item Friendly error messages
\end{itemize}

\section{Diagram: Error Handling Flow}

\begin{center}
\begin{tikzpicture}[node distance=2.5cm]
\node (try) [draw, rectangle, fill=yellow!20, minimum height=1cm] {Try Block};
\node (error) [draw, diamond, aspect=2, fill=green!20, below of=try, yshift=-0.5cm] {Error?};
\node (except) [draw, rectangle, fill=red!20, right of=error, xshift=2cm, minimum height=1cm] {Except Block};
\node (else) [draw, rectangle, fill=blue!20, left of=error, xshift=-2cm, minimum height=1cm] {Else Block};

\draw[->] (try) -- (error);
\draw[->] (error) -- node[above]{Yes} (except);
\draw[->] (error) -- node[above]{No} (else);
\end{tikzpicture}
\end{center}

\section{Best Practices}

\begin{itemize}
    \item Handle specific errors
    \item Use clear error messages
    \item Avoid empty except blocks
\end{itemize}

\section{Key Concepts}

\begin{itemize}
    \item Errors stop programs
    \item \texttt{try-except} handles errors
    \item \texttt{else} runs when no error
    \item \texttt{finally} always runs
\end{itemize}

\section{Common Beginner Mistakes}

\begin{itemize}
    \item Using except without knowing error
    \item Hiding errors silently
    \item Ignoring error messages
\end{itemize}

\section{End-of-Chapter Exercises}

\subsection{Very Short Questions}

\begin{enumerate}
    \item What is an error?
    \item Why is error handling important?
    \item Which block always runs?
\end{enumerate}

\subsection{Short Answer Questions}

\begin{enumerate}
    \item Explain try-except using a real-life example.
    \item Difference between else and finally.
    \item Why should we avoid empty except?
\end{enumerate}

\subsection{Activity-Based Exercises}

\begin{enumerate}
    \item Create a safe division program.
    \item Handle missing file error.
    \item Write a program that accepts only positive numbers.
\end{enumerate}

\section{Teacher Notes}

\marginnote{
\textbf{Learning Objectives}
\begin{itemize}
    \item Students understand runtime errors
    \item Students use try-except blocks
    \item Students build crash-free programs
\end{itemize}

\textbf{Teaching Tips}
\begin{itemize}
    \item Compare with exam mistakes
    \item Demonstrate crash vs safe program
    \item Encourage user-friendly messages
\end{itemize}

\textbf{Assessment}
\begin{itemize}
    \item Correct exception handling
    \item Clear error messages
\end{itemize}
}